% generator_I01.tex -- Prop I.1: construct an equilateral triangle on a given line.
% Reference: jemmybutton vector recreation, page 25.
% Build: lualatex generator_I01.tex  ->  generator_I01.pdf
\documentclass[12pt]{article}
\usepackage{geometry}
\geometry{a5paper, margin=1.5cm}
\usepackage[lining]{ebgaramond}
\usepackage{amssymb}
\usepackage{byrne}
\def\mpPre{u := 1cm; textLabels := false;}

\begin{document}

% FIGURE (defined once; all proof steps reuse these points)
\defineNewPicture[1/2]{%
    pair A, B, C;%
    path P[];%
    numeric r;%
    r := 3/2u;%
    A := (0, 0);%
    B := (r, 0);%
    P1 := fullcircle scaled 2r;%
    P2 := fullcircle scaled 2r shifted B;%
    C := P1 intersectionpoint P2;%
    % Layer-1 assertions: all sides equal, C lies on both construction circles
    if abs(abs(A-B) - abs(B-C)) > 0.01: errmessage "ASSERT: AB != BC"; fi;%
    if abs(abs(B-C) - abs(C-A)) > 0.01: errmessage "ASSERT: BC != CA"; fi;%
    if abs(abs(C-A) - r) > 0.01: errmessage "ASSERT: C not on circle A"; fi;%
    if abs(abs(C-B) - r) > 0.01: errmessage "ASSERT: C not on circle B"; fi;%
    byLineDefine(A, B, byblack, SOLID_LINE, REGULAR_WIDTH);%
    byLineDefine(B, C, byred,   SOLID_LINE, REGULAR_WIDTH);%
    byLineDefine(C, A, byyellow,SOLID_LINE, REGULAR_WIDTH);%
    draw byNamedLineSeq(-1)(AB,CA,BC);%
    draw byCircle.A(A, B, byblue, 0, 0,  1/2);%
    draw byCircle.B(B, A, byred,  0, 0, -1/2);%
    draw byLabelsOnPolygon(A, C, B)(ALL_LABELS, 1);%
}

% STATEMENT
\begin{center}
\textbf{Book~I.\enspace Prop.\ I. Prob.}
\end{center}

\vspace{0.4cm}
\noindent
\begin{minipage}[t]{0.50\linewidth}
\itshape On a given finite straight line
(\drawUnitLine{AB}) to describe an equilateral triangle.
\end{minipage}%
\hfill
\begin{minipage}[t]{0.46\linewidth}
\centering
\drawCurrentPicture
\end{minipage}

\vspace{0.6cm}

% PROOF
\begin{center}
Describe \offsetPicture{15pt}{0pt}{\drawFromCurrentPicture{%
    draw byNamedLine(AB);%
    draw byNamedCircle(A);%
    draw byLabelLineEnd(A, B, 0);%
    draw byLabelLineEnd(B, A, 1);%
}} and \offsetPicture{15pt}{0pt}{\drawFromCurrentPicture{%
    draw byNamedLine(AB);%
    draw byNamedCircle(B);%
    draw byLabelLineEnd(A, B, 1);%
    draw byLabelLineEnd(B, A, 0);%
}} (post.\ III);\\[4pt]
draw \drawUnitLine{CA} and \drawUnitLine{BC} (post.\ I).\\[4pt]
Then will \drawLine[bottom][triangleABC]{AB,CA,BC} be equilateral.\\[8pt]
For \quad $\drawUnitLine{AB} = \drawUnitLine{CA}$ \quad (def.\ 15);\\
and \quad $\drawUnitLine{AB} = \drawUnitLine{BC}$ \quad (def.\ 15);\\
$\therefore \drawUnitLine{CA} = \drawUnitLine{BC}$ \quad (ax.\ I);\\[6pt]
and therefore \triangleABC\ is the equilateral triangle required.
\end{center}

\begin{flushright}Q.\ E.\ D.\end{flushright}

\end{document}
